
\section{Target-control records}
\label{app:targetrecords}

The main-text examples use the saved outputs
\texttt{runs151/R11\_r1.png}, \texttt{runs151/R13\_r3.png}, and
\texttt{runs152/K34b\_r4.png}. The observations in
Table~\ref{tab:targetoutputs} describe the images below. They retain the
visible rendering style and body configuration as separate properties.
Geometric chest marks are recorded explicitly.

\begin{figure*}[t]
\centering
\includegraphics[width=0.31\textwidth]{figures/target_standing.png}\hfill
\includegraphics[width=0.31\textwidth]{figures/target_crouching.png}\hfill
\includegraphics[width=0.31\textwidth]{figures/target_rear.png}
\caption{\textbf{Original target-control outputs.} Left: batch 151, R11,
repeat 1. Centre: batch 151, R13, repeat 3. Right: batch 152, K34b,
repeat 4. These are the original images used for the attribute observations
in Table~\ref{tab:targetoutputs}.}
\label{fig:targetrecords}
\end{figure*}

Batch 156 contains 39 saved request records. Groups M, N, P, and G each
contain three pose settings with three attempts per setting. Group Q contains
three further pose-wording attempts. N11 and P11 each return one image;
the other groups return none. Table~\ref{tab:rendercontext} uses groups
M, N, P, and G. The source filenames and record hashes are frozen in
\texttt{restructure\_evidence.json}. These records supplement the earlier
corpus census rather than changing its frozen totals.

\section{Supplementary component ablations}
\label{app:ablationdetail}

The historical component sweeps vary the opening context, texture depth,
and carrier form separately. Their original summary is reproduced in
Figure~\ref{fig:suppcomponents}. These measurements concern returned images;
target-attribute attainment requires the separate checks in
\S\ref{sec:targetmatch}.

\begin{figure}[t]
\centering
\includegraphics[width=\linewidth]{figures/figD_three_arms.png}
\caption{\textbf{Component sweeps.} The panels vary the opening context,
texture-layer count, and carrier form. Each panel represents its own batch;
the layer sweep covers one through ten layers.}
\label{fig:suppcomponents}
\end{figure}
